%*******************************************************
% Résumé
%*******************************************************
\pdfbookmark[1]{Résumé}{resume}

\chapter*{Résumé}

La transition numérique reconfigure la relation entre systèmes économiques et infrastructures énergétiques. Pourtant, les cadres analytiques mobilisés pour évaluer cette transformation restent fragmentés~: les modèles économiques traitent l'ICT comme un choc de productivité exogène, tandis que les modèles énergie-climat peinent à endogénéiser la demande matérielle croissante du calcul. Cette thèse examine les dimensions économiques et énergétiques de l'infrastructure ICT à travers une approche pluridisciplinaire combinant histoire de l'innovation, analyse économique structurelle et modélisation énergie-climat.

Le chapitre~1 retrace l'histoire phénoménologique de l'ICT du télégraphe au datacenter, en montrant comment les questions économiques et énergétiques émergent naturellement de la trajectoire technique. Le chapitre~2 développe une analyse économique systémique de l'hypothèse selon laquelle le calcul constitue un nouveau régime structurel. Le chapitre~3 confronte les promesses de la décarbonation par l'ICT aux réalités empiriques~: barrières, échecs d'adoption, conditions systémiques. Les chapitres~4 à~6 mettent en \oe uvre cette analyse dans le cadre de modélisation IMACLIM, en construisant des scénarios qui endogénéisent le couplage ICT-énergie. Le chapitre~7 discute les implications en termes de politiques publiques.

L'argument central est que l'infrastructure ICT fonctionne comme un \emph{pharmakon}~--- simultanément productrice et réductrice d'entropie~--- et que le datacenter, lieu de convergence du GPT-Électricité et du GPT-ICT, représente une rupture structurelle exigeant de nouveaux outils analytiques.
